\section[Geometry and architecture]{Geometry conventions and engine architecture (LEA/TOP, ARCH)}
\label{sec:geometry-arch}
% MOC -> MDO data flow (current default vs proposed M0a/M0b). Drawn by writer_t.
% Sources: plan_record avenues.ARCH, planner_findings, milestone M0; raosim/mdo/engine.py (solve_engine, engine_residual).
\begin{figure}[htbp]
\centering
\begin{tikzpicture}[x=1in,y=1in,
  box/.style={draw=teal,line width=0.5pt,fill=white,rounded corners=2pt,text=ink,
    font=\sffamily\scriptsize,align=left,inner xsep=3pt,inner ysep=2pt},
  boxA/.style={box,draw=teal,dash pattern=on 2pt off 1.2pt},
  boxB/.style={box,draw=accent,dash pattern=on 2pt off 1.2pt},
  boxR/.style={box,draw=muted,text=muted,dash pattern=on 1pt off 1pt},
  cur/.style={draw=ink,line width=0.5pt,rounded corners=1.6pt,-{Stealth[length=1.6mm,width=1.1mm]}},
  ma/.style={draw=teal,line width=0.6pt,rounded corners=1.6pt,-{Stealth[length=1.6mm,width=1.1mm]}},
  mb/.style={draw=accent,line width=0.6pt,rounded corners=1.6pt,-{Stealth[length=1.6mm,width=1.1mm]}},
  rep/.style={draw=muted,line width=0.5pt,dash pattern=on 2pt off 1.2pt,-{Stealth[length=1.6mm,width=1.1mm]}},
  lab/.style={font=\sffamily\scriptsize,inner sep=1pt,align=center}]
% ---- band: the per-design-point sizing loop (raosim/mdo/engine.py solve_engine) ----
\fill[pale,rounded corners=3pt] (0,0.55) rectangle (6.8,-1.34);
\node[font=\sffamily\scriptsize\bfseries,text=teal,anchor=west,inner sep=0pt] at (0.05,0.465) {Sizing loop};
\node[font=\sffamily\scriptsize\bfseries,text=muted,anchor=west,align=left,inner sep=0pt] at (0.05,0.86)
  {Outside the\\sizing loop};
% ---- outside the loop ----
\node[boxA,text width=1.25in-6pt,minimum height=0.46in] (tab) at (2.00,0.86)
  {\textbf{M0a tables}\\ $\lambda=C_{F,\mathrm{MOC}}/C_{F,1D}$\\ $p_{w,e}$; $(\varepsilon,L\%,\gamma,R_d)$};
\node[boxA,text width=1.0in-6pt,minimum height=0.46in] (swp) at (3.64,0.86)
  {\textbf{Offline sweeps}\\ N0 forward net\\ approval: D10};
\node[boxR,text width=1.1in-6pt,minimum height=0.46in] (post) at (4.95,0.86)
  {\textbf{MOC post-solve}\\ Rao BVP or\\ forward net};
\node[font=\sffamily\scriptsize,text=muted,anchor=west,align=left] at (5.55,0.92) {comparison and\\report only; nothing\\returns to sizing};
% ---- inside the loop: main row ----
\node[box,text width=0.85in-6pt,minimum height=0.74in] (in) at (0.425,0)
  {\textbf{Design point}\\ $P_c$, $\varepsilon$, O/F\\ $F$, $p_a$\\ $\gamma$, $c^*$, $T_c$};
\node[box,text width=1.3in-6pt,minimum height=0.74in] (cf) at (2.00,0)
  {\textbf{Thrust coefficient}\\
   now: $C_F=\eta_{CF}\,C_{F,1D}$\\
   $\eta_{CF}$ fixed, 0.98--0.99\\
   {\color{teal}M0a: $C_F=\lambda\,C_{F,1D}$}\\
   {\color{accent}M0b: $C_F=C_{F,\mathrm{MOC}}$}};
\node[box,text width=1.15in-6pt,minimum height=0.74in] (res) at (3.525,0)
  {\textbf{Sizing residuals}\\ $R_1=F-C_F P_c A_t$\\ $R_2=\dot m-P_c A_t/c^*$\\ Newton, IFT\\ $\to R_t,\ \dot m$};
\node[box,text width=1.1in-6pt,minimum height=0.74in] (geo) at (4.95,0)
  {\textbf{Contour provider}\\ TOP/B\'ezier default\\ $r(x)$, station grid\\ convention: G0};
\node[box,text width=0.9in-6pt,minimum height=0.46in] (th) at (6.35,0)
  {\textbf{Thermal}\\ Bartz heat flux,\\ regen cooling};
% ---- inside the loop: lower row ----
\node[box,text width=1.05in-6pt,minimum height=0.6in] (sep) at (1.875,-1.00)
  {\textbf{Separation}\\ now: Schmucker,\\ Q1D exit state\\ opt-in: wall, V1};
\node[boxB,text width=1.2in-6pt,minimum height=0.6in] (fn) at (3.66,-1.00)
  {\textbf{M0b forward net}\\ N0 primal, N1 totals\\ gas mode, geometry\\ typed fallback};
\node[box,text width=0.9in-6pt,minimum height=0.46in] (st) at (6.35,-1.00)
  {\textbf{Structure}\\ wall loads, mass};
% ---- current default (ink) ----
\draw[cur] (in.east) -- node[lab,above,text=ink,pos=0.34]{$\varepsilon,\gamma$\\$P_c,p_a$} (cf.west);
\draw[cur] (cf.east) -- node[lab,above,text=ink]{$C_F$} (res.west);
\draw[cur] (res.east) -- node[lab,above,text=ink]{$R_t$} (geo.west);
\draw[cur] ([yshift=0.08in]geo.east) -- node[lab,above,text=ink,pos=0.40]{$r(x)$} ([yshift=0.08in]th.west);
\draw[cur] ([yshift=0.08in]geo.east) -- (5.78,0.08) -- (5.78,-0.92) -- ([yshift=0.08in]st.west);
\draw[cur] (1.62,0 |- cf.south) -- node[lab,right,text=ink,align=left]{Q1D exit\\$M_e,\ p_e$} (1.62,0 |- sep.north);
\draw[rep] (geo.north) -- node[lab,right,text=muted]{contour} (post.south);
% ---- M0a (teal) ----
\draw[ma] (swp.west) -- node[lab,above,text=teal]{$C_{F,\mathrm{MOC}}$} (tab.east);
\draw[ma] (2.25,0 |- tab.south) -- node[lab,right,text=teal]{$\lambda$ replaces $\eta_{CF}$} (2.25,0 |- cf.north);
\draw[ma] (tab.west) -- (1.27,0.86) -- (1.27,-1.00) -- (sep.west);
\node[lab,text=teal,anchor=east] at (1.25,-0.78) {$p_{w,e}$};
% ---- M0b (accent) ----
\draw[mb] (4.72,0 |- geo.south) -- (4.72,-0.50) -- (3.98,-0.50) -- (3.98,0 |- fn.north);
\node[lab,text=accent,anchor=north west] at (4.30,-0.53) {wall $r(x)$, $R_d$, $R_u$};
\draw[mb] (3.30,0 |- fn.north) -- (3.30,-0.58) -- (2.52,-0.58) -- (2.52,0 |- cf.south);
\node[lab,text=accent,anchor=north,align=center] at (2.76,-0.60) {$C_{F,\mathrm{MOC}}$,\\N1 totals};
\draw[mb] (fn.west) -- node[lab,below,text=accent,align=center]{wall\\$p$, $M$} (sep.east);
\coordinate (fnE) at ([yshift=-0.08in]fn.east);
\draw[mb] (fnE) -- node[lab,below,text=accent,pos=0.47]{wall $p$, $T$, $M$; $\delta^*$ (V0)} (5.66,0 |- fnE) -- ([yshift=-0.08in]st.west);
\draw[mb] (5.66,0 |- fnE) -- (5.66,-0.08) -- ([yshift=-0.08in]th.west);
% ---- legend ----
\begin{scope}[font=\sffamily\scriptsize]
  \draw[cur] (0.30,-1.52) -- (0.62,-1.52); \node[anchor=west,text=ink] at (0.65,-1.52) {current default (Q1D closure)};
  \draw[ma]  (2.30,-1.52) -- (2.62,-1.52); \node[anchor=west,text=ink] at (2.65,-1.52) {M0a, opt-in};
  \draw[mb]  (3.55,-1.52) -- (3.87,-1.52); \node[anchor=west,text=ink] at (3.90,-1.52) {M0b, opt-in after N1};
  \draw[rep] (5.30,-1.52) -- (5.62,-1.52); \node[anchor=west,text=ink] at (5.65,-1.52) {report only};
\end{scope}
\end{tikzpicture}
\caption{Data flow from the MOC to the MDO. In the current default the loop closes thrust and throat radius with
$C_F=\eta_{CF}C_{F,1D}$ and evaluates the Schmucker margin at the Q1D exit state before requesting the TOP/B\'ezier contour,
so the closure is blind to bell length and wall shape; the MOC is a post-solve check only. M0a (teal) replaces $\eta_{CF}$ by
$\lambda=C_{F,\mathrm{MOC}}/C_{F,1D}$ tabulated offline from forward-net sweeps (E04: 0.9672 for Nozzle-B); M0b (orange)
moves the forward net into the loop after N1, feeding $C_F$ and wall states to the residual and the consumers.}
\label{fig:moc_mdo_flow}
\end{figure}

\textbf{SP-125 length.} SP-125 measures bell length against a $15^\circ$ cone with
the same throat area, radius below the throat and area ratio
\cite[PDF p.~100, printed p.~91]{huzel1967}:
\begin{equation}
L_n=\frac{R_t(\sqrt{\epsilon}-1)+R(\sec\alpha-1)}{\tan\alpha}
\end{equation}
\cite[PDF p.~99, printed p.~90, Eq.~4-7]{huzel1967}, with $R=0.382R_t$ in the worked
example (PDF p.~107, printed p.~98).

\textbf{Diagnosis (E02).} The working-tree TOP/B\'ezier contour now uses this rounded
cone. The LEA case, however, keeps \code{length_pct} $=1500/L_{\mathrm{sharp}}$,
derived from the sharp-corner cone, so the host gives $L/R_t=15.044818739$ against
$15\pm0.02$.

Of 20 $\tan15^\circ$ sites (12 in \code{raosim/}, 8 in tests), 19 use the sharp corner.
In \code{raosim/} 11 of 12 do, including the MDO grid, the engine collapse screen, the
Rao BVP and the optimiser. \code{tests/test_nozzle_geometry.py} has four working-tree
failures: one from the convention and three from chart changes.

\textbf{The source reference.} Schomberg et al.'s Table~3 gives $x_{mn}=15.0r_t$ without
defining it. Their Fig.~5 suggests it is measured from the end of a $0.5y_t$ throat arc
\cite[PDF pp.~5--6 of the held EUCASS 2015 copy]{schomberg2019}; this is an inference.
The restated geometry ($R_d=0.5R_t$, $x_E-x_N=15.0$) is $15.2796R_t$ from throat to
exit.

Its contour fit is better in the throat-origin alignment (RMS 0.0714 against 0.0870,
E02b). With $x$ measured from N the two are almost equal (0.148 against 0.152), so this
evidence is weak. The $15\pm0.02$ check was circular (MOC-30).

\textbf{Thrust closure (ARCH).} The engine residual uses
$C_F=\eta_{C_F}C_{F,\mathrm{1D}}$ and the margin $P_e/p_{\mathrm{sep}}(M_e)-1.2$ from
the Q1D exit state, so it is blind to bell length and shape (Fig.~\ref{fig:moc_mdo_flow}).

E04 gives $\lambda=C_{F,\mathrm{MOC}}/C_{F,\mathrm{1D}}=0.9672$ for Nozzle-B with the
production start line ($C_D=1.0075$), and about 0.955 in the KL-V3 replay. The MDO
instead uses fixed $\eta_{C_F}$ values of 0.98--0.99. The MOC wall exit pressure is
$0.0465p_0$ against $0.0302p_0$ from Q1D.

\begingroup\footnotesize
\par\addvspace{\LTpre}{\centering\begin{tabular}{@{}P{0.30\linewidth}P{0.50\linewidth}P{0.13\linewidth}@{}}
\toprule\textbf{Approach} & \textbf{Assessment} & \textbf{Verdict}\\\midrule
Explicit \code{length_convention} helper at every site, defaults kept & Removes the hidden split; JAX-safe & \textbf{G0}\\
Restate benchmark cases in the source's quantities & Needs D11 for LEA & with G0\\
Revert host to sharp corner, or switch all to SP-125 & Silent default change & D7 only\\
Keep Q1D thrust closure & Length-blind & default\\
MOC correction table $\lambda(\epsilon,L\%,\gamma,R_d)$ (M0a) & Cheap, smooth; must record its start-line method & \textbf{M0a}\\
Forward net in the residual (M0b) & Consistent; needs N1 & \textbf{M0b}\\
Rao BVP in the inner loop & About 25 min per accepted pair & reject\\
\bottomrule
\end{tabular}\par}\addvspace{\LTpost}
\endgroup

\begingroup\footnotesize
\par\addvspace{\LTpre}{\centering\begin{tabular}{@{}P{0.09\linewidth}P{0.36\linewidth}P{0.30\linewidth}P{0.17\linewidth}@{}}
\toprule\textbf{ID} & \textbf{Numerical verification} & \textbf{Benchmark agreement} & \textbf{Physical validation}\\\midrule
C0 & Contract round trips; an invalid provider cannot yield accepted performance & n/a & n/a\\
G0 & Both lengths to $10^{-9}$ (done); per-site and JAX-parity tests & LEA and Vulcain S1 geometry pass unchanged tolerances on the source quantity (after D11) & n/a\\
M0 & Perturbations reach the right residuals; table holdout error; Taylor tests (M0b) & n/a & none\\
D0 & Same-context CLI/API/MDO replay & as A0 & as A0\\
\bottomrule
\end{tabular}\par}\addvspace{\LTpost}
\endgroup
